% Fig — the wrap-step bug: a lunation grid spans [0, P-dt], so a march that
% stops at t[n-1] and relabels it t=0 slips phase by dt every cycle.
% Mirrors solver.solve_pixel (fixed 2026-07-03; audit F-5 addendum III).
% Body only; the book supplies \caption/\label.
\begin{tikzpicture}[
  font=\footnotesize,
  every node/.style={align=center},
  hop/.style={-{Stealth[length=1.8mm]}, line width=0.8pt, draw=dim, line cap=round},
  hopbad/.style={hop, draw=coral, dashed},
  hopfix/.style={hop, draw=forest},
  lbl/.style={font=\scriptsize\color{dim}},
  tag/.style={rectangle, rounded corners=2pt, line width=0.9pt, inner sep=4pt,
              text width=44mm, font=\scriptsize}]

  % ---- top strip: the buggy march --------------------------------------
  \draw[line width=0.9pt, color=char] (0,0) -- (8.4,0);
  \foreach \x in {0,1.2,2.4,3.6,4.8,6.0} \draw[color=char] (\x,-0.07) -- (\x,0.07);
  \draw[color=char] (7.2,-0.07) -- (7.2,0.07);
  \draw[color=char, line width=1.2pt] (8.4,-0.10) -- (8.4,0.10);
  \node[lbl, below=2pt] at (0,0)   {$t_0{=}0$};
  \node[lbl, below=2pt] at (3.6,0) {$\cdots$};
  \node[lbl, below=2pt] at (7.2,0) {$t_{n-1}{=}P{-}\Delta t$};
  \node[lbl, below=2pt] at (8.4,0) {$P$};
  \foreach \a/\b in {0/1.2, 1.2/2.4, 2.4/3.6, 3.6/4.8, 4.8/6.0, 6.0/7.2}
    \draw[hop] (\a,0.12) to[bend left=45] (\b,0.12);
  \draw[hopbad] (7.2,0.12) to[bend left=45] (8.4,0.12);
  \node[lbl, color=coral] at (7.8,0.78) {never taken};
  % the relabel: the P-dt state jumps back and is called t=0 (bows DOWN,
  % well clear of the axis, its ticks, and the hops)
  \draw[hopbad, dash pattern=on 2.5pt off 2pt]
    (7.2,-0.45) to[bend left=16]
    node[lbl, below=3pt, color=coral, pos=0.5]
    {state at $P{-}\Delta t$ relabelled as $t{=}0$ of the next cycle}
    (0,-0.45);
  \node[tag, draw=coral, fill=coral!7, anchor=west] at (9.0,-0.2)
    {\textbf{buggy march}: the last step is never taken --- the cycle
     advances only $P{-}\Delta t$, a phase slip of $\Delta t$ per lunation
     $\Rightarrow$ even the converged state is $O(\Delta t)$ wrong};

  % ---- bottom strip: the fixed march ------------------------------------
  \begin{scope}[yshift=-2.7cm]
    \draw[line width=0.9pt, color=char] (0,0) -- (8.4,0);
    \foreach \x in {0,1.2,2.4,3.6,4.8,6.0} \draw[color=char] (\x,-0.07) -- (\x,0.07);
    \draw[color=char] (7.2,-0.07) -- (7.2,0.07);
    \draw[color=char, line width=1.2pt] (8.4,-0.10) -- (8.4,0.10);
    \node[lbl, below=2pt] at (0,0)   {$t_0{=}0$};
    \node[lbl, below=2pt] at (3.6,0) {$\cdots$};
    \node[lbl, below=2pt] at (7.2,0) {$P{-}\Delta t$};
    \node[lbl, below=2pt] at (8.4,0) {$P\equiv 0$};
    \foreach \a/\b in {0/1.2, 1.2/2.4, 2.4/3.6, 3.6/4.8, 4.8/6.0, 6.0/7.2}
      \draw[hop] (\a,0.12) to[bend left=45] (\b,0.12);
    \draw[hopfix, line width=1.1pt] (7.2,0.12) to[bend left=45]
      node[lbl, above=2pt, xshift=-9mm, color=forest]
      {the wrap step ($S = S(0)$)} (8.4,0.12);
    \node[tag, draw=forest, fill=forest!10, anchor=west] at (9.0,0.1)
      {\textbf{fixed march}: one more step closes the cycle exactly; its
       result \emph{is} the next cycle's $t{=}0$ state (and surface
       temperature) $\Rightarrow$ the attractor is second-order};
  \end{scope}
\end{tikzpicture}
